\documentclass[11pt,a4paper]{article}
\usepackage{amsmath,amssymb,amsthm}
\usepackage{geometry}
\geometry{margin=2.5cm}
\usepackage{hyperref}
\usepackage{booktabs}
\usepackage{authblk}

\newtheorem{theorem}{Theorem}
\newtheorem{proposition}[theorem]{Proposition}
\newtheorem{corollary}[theorem]{Corollary}
\newtheorem{definition}[theorem]{Definition}
\newtheorem{assumption}[theorem]{Assumption}
\newtheorem{remark}[theorem]{Remark}

\newcommand{\Hthree}{\mathbb{H}^3}

\begin{document}

\title{Discrete Mixing Operators from Boundary Sector Geometry}
\author[1]{Marvin L. Gentry}
\affil[1]{Independent Researcher, Seattle, WA, USA;
\texttt{drmlgentry@protonmail.com};
ORCID: 0009-0006-4550-2663}
\date{\today}
\maketitle

\begin{abstract}
The CKM and PMNS mixing matrices of the Standard Model lack a
first-principles geometric origin.
Let $F$ be a finite set of fermion labels; each $f\in F$ is assigned
a distinct boundary direction $\xi(f)$ on the sphere at infinity
$S^2=\partial\mathbb{H}^3$ of a compact hyperbolic 3-manifold.
Using a Hilbert space of boundary data with a continuous family of
normalized vectors $|\xi\rangle$, we define sector subspaces
$H_S=\operatorname{span}\{|\xi(f)\rangle:f\in S\}$ for any ordered
subset $S\subset F$.
The mixing matrix $U_{ST}$ between two sectors $S,T$ of equal size
is the matrix of the orthogonal projection $P_T|_{H_S}$ in the
Gram-Schmidt orthonormal bases of $H_S$ and $H_T$.
We prove that under the discrete assignment (the set
$\Xi=\{\xi(f):f\in F\}$ is finite and locally rigid) the accessible
mixing matrices form a finite set modulo
$U(1)^n\times U(1)^n$ rephasings.
Consequently, no continuous deformation of the boundary points can
produce continuous drift in any rephasing-invariant mixing observable
--- a structural constraint that any realistic model must satisfy.
In companion papers~\cite{Gentry:CKM,Gentry:PMNS}, this framework
is applied to two arithmetic hyperbolic 3-manifolds, reproducing the
CKM and PMNS matrices with Frobenius fitness $0.016482$ and
$0.005087$ respectively.
\end{abstract}

\section{Introduction}

Mixing matrices, such as the CKM matrix in the quark sector and the
PMNS matrix in the lepton sector, encode the mismatch between flavor
and mass eigenstates in the Standard Model~\cite{PDG}.
Despite their fundamental importance, their origin remains unexplained
by the gauge structure alone.
Many attempts invoke discrete flavor symmetries, extra dimensions, or
string-inspired constructions~\cite{AltarelliFeruglio,KingLuhn}.
In all such approaches, the mixing matrix ultimately measures the
overlap between states associated with different sectors of the theory.

In this paper we take a kinematic, model-independent perspective.
We show that if fermions are assigned to distinct boundary directions
of a compact hyperbolic 3-manifold, the mixing matrices arising from
the overlap of boundary-localized states are determined solely by the
discrete boundary data, and form a discrete set modulo sector-local
rephasings.
A corollary is that no continuous deformation of physical parameters
can produce continuous drift in any rephasing-invariant mixing observable.

The construction is deliberately minimal: a Hilbert space of boundary
data, a family of boundary-localized states, and a discrete assignment
of fermions to boundary directions.
No dynamics, no Lagrangian, and no ultraviolet completion are assumed.

In companion papers~\cite{Gentry:CKM,Gentry:PMNS}, this framework is
applied to two arithmetic hyperbolic 3-manifolds:
$M_{\mathrm{PMNS}}=m003(-2,3)$ (OrientableClosedCensus[1]) and
$M_{\mathrm{CKM}}=m006(-5,2)$ (OrientableClosedCensus[43]).
The Borel (N-factor Iwasawa) construction on
$M_{\mathrm{PMNS}}$ reproduces the PMNS matrix with Frobenius fitness
$0.005087$ against PDG~2024 values --- the global minimum of the
construction ($p<10^{-4}$ vs $50{,}000$ Haar-random unitaries).
The K-factor construction on $M_{\mathrm{CKM}}$ gives fitness
$0.016482$ with Cabibbo angle error $0.19\%$.
The present paper provides the kinematic guarantee that these results
are not accidents of continuous parameter tuning.

\section{Geometric Setup}
\label{sec:setup}

Let $\Gamma$ be a torsion-free discrete group acting properly
discontinuously and cocompactly on $\mathbb{H}^3$ by
orientation-preserving isometries.
The quotient $M=\mathbb{H}^3/\Gamma$ is a compact oriented
hyperbolic 3-manifold with boundary at infinity $S^2=\partial\mathbb{H}^3$.

\begin{definition}[Boundary direction]
A \emph{boundary direction} is a point $\xi\in S^2$.
For $\gamma\in\Gamma$, the associated boundary direction is the ideal
endpoint of the geodesic ray from a fixed basepoint $o\in\mathbb{H}^3$
through $\gamma\cdot o$.
\end{definition}

\begin{assumption}[Discrete sector assignment]
\label{ass:discrete}
Each fermion flavor $f\in F$ is assigned a unique boundary direction
$\xi(f)\in S^2$ by a deterministic, discrete selection rule: the image
$\Xi=\{\xi(f):f\in F\}$ is finite, and the assignment is locally rigid.
\end{assumption}

\begin{remark}[Geometric realization]
Fix $o\in\mathbb{H}^3$. For each $f$, let $\gamma_f\in\Gamma$ be the
unique element minimizing $d(o,\gamma\cdot o)$ in the assignment class
of $f$ (uniqueness follows from proper discontinuity and
torsion-freeness~\cite{BridsonHaefliger}).
The boundary direction $\xi(f)$ is the ideal endpoint of the geodesic
from $o$ through $\gamma_f\cdot o$.
In the HFG application~\cite{Gentry:CKM,Gentry:PMNS},
$\gamma_f$ is a short word in $\pi_1(M)$ and $\xi(f)=\hat{n}(\gamma_f)$
is the geodesic axis extracted via the Pauli decomposition of
$\log\rho(\gamma_f)$.
\end{remark}

\section{Boundary Hilbert Space and Sector Subspaces}
\label{sec:boundary}

Let $\mathcal{H}$ be a separable complex Hilbert space (concretely $L^2(S^2)$).

\begin{definition}[Boundary-localized state]
For each $\xi\in S^2$, fix a normalized vector $|\xi\rangle\in\mathcal{H}$
with $\xi\mapsto|\xi\rangle$ continuous, $\langle\xi|\xi\rangle=1$,
and $|\langle\xi|\zeta\rangle|<1$ for $\xi\neq\zeta$.
An explicit family is given by the Poisson kernel
$|\xi\rangle\propto P(\cdot,\xi)$ normalized in $L^2(S^2)$~\cite{Ratcliffe}.
\end{definition}

\begin{definition}[Sector subspace and preferred basis]
A \emph{sector} $S$ is a finite ordered subset of $F$.
The \emph{sector subspace} is
$H_S=\operatorname{span}\{|\xi(f)\rangle:f\in S\}\subset\mathcal{H}$,
with $\dim H_S=|S|$.
The \emph{preferred ONB} $\{|S,1\rangle,\ldots,|S,n\rangle\}$ of $H_S$
is the Gram-Schmidt orthonormalization of
$(|\xi(f_1)\rangle,\ldots,|\xi(f_n)\rangle)$
with all pivot coefficients real and positive,
removing the $U(1)^n$ phase ambiguity.
\end{definition}

\section{Mixing Matrices as Overlap Operators}
\label{sec:mixing}

\begin{definition}[Mixing matrix]
For sectors $S,T$ of equal cardinality $n$, the \emph{mixing matrix}
$U_{ST}$ has entries $(U_{ST})_{ji}=\langle T,j|S,i\rangle$.
\end{definition}

\begin{proposition}[Unitarity and geometric dependence]
When $H_S=H_T$, $U_{ST}\in U(n)$; in general it is a partial isometry.
$U_{ST}$ is determined entirely by $\Xi_{ST}=\{\xi(f):f\in S\cup T\}$
and the inner product on $\mathcal{H}$: all Gram-Schmidt coefficients
are algebraic expressions in $\langle\xi(f)|\xi(g)\rangle$,
$f,g\in S\cup T$.
\end{proposition}

\section{Rephasing Invariance}
\label{sec:rephasing}

Physical observables are invariant under $U_{ST}\to D_T U_{ST}D_S^\dagger$,
$D_S,D_T\in U(1)^n$.

\begin{proposition}[Rephasing-invariant quantities]
Invariant under sector-local rephasings:
(i)~moduli $|(U_{ST})_{ji}|^2$;
(ii)~quartic products
$J_{ji,j'i'}=(U_{ST})_{ji}(U_{ST})_{j'i'}\overline{(U_{ST})_{ji'}}\,
\overline{(U_{ST})_{j'i}}$,
generalizing the Jarlskog invariant~\cite{Jarlskog};
(iii)~$\det U_{ST}$ up to an overall phase.
For $n=3$, $\operatorname{Im}J_{ji,j'i'}$ reduces to $J_{\mathrm{CP}}$.
\end{proposition}

\section{Stability and Discreteness}
\label{sec:stability}

\begin{theorem}[Stability]
\label{thm:stability}
Let $\delta_{\min}=\min_{f\neq f'}d_{S^2}(\xi(f),\xi(f'))$ and
$\varepsilon_0=\delta_{\min}/2$.
For any deformation of magnitude $\varepsilon\leq\varepsilon_0$,
the perturbed configuration determines the same preferred bases and
the same $U_{ST}$.
\end{theorem}

\begin{proof}
For $\varepsilon\leq\varepsilon_0$, each $\xi'(f)$ remains closer to
$\xi(f)$ than to any $\xi(g)$, $g\neq f$, so labels are preserved
and linear independence holds.
Gram-Schmidt coefficients depend continuously on inner products;
since the discrete assignment rule (Assumption~\ref{ass:discrete})
is unchanged, $U_{ST}$ depends only on the unchanged configuration
class (by the geometric dependence proposition).
\end{proof}

\begin{theorem}[Discreteness of mixing classes]
\label{thm:discreteness}
Under Assumption~\ref{ass:discrete}, the set
$\mathcal{M}_{ST}=\{U_{ST}:\Xi_{ST}\text{ admissible}\}$
is finite, hence discrete, in
$U(n)/(U(1)^n\times U(1)^n)$.
\end{theorem}

\begin{proof}
The set of admissible $\Xi_{ST}$ is finite.
The map $\Xi_{ST}\mapsto U_{ST}$ is continuous, hence sends a finite
set to a finite set.
Quotienting by the rephasing action preserves finiteness.
\end{proof}

\begin{corollary}[No continuous drift]
For deformations of magnitude $\varepsilon\leq\varepsilon_0$, no
continuous drift occurs in any rephasing-invariant observable.
For larger deformations, the mixing matrix changes only by a discrete
jump within $\mathcal{M}_{ST}$.
\end{corollary}

\begin{proof}
Rephasing-invariant observables are continuous on
$U(n)/(U(1)^n\times U(1)^n)$;
their restriction to the finite set $\mathcal{M}_{ST}$ takes only
finitely many values.
\end{proof}

\section{Connection to Hyperbolic Flavor Geometry}
\label{sec:hfg}

The abstract framework is concretely realized in the HFG
programme~\cite{Gentry:CKM,Gentry:PMNS,Gentry:Unified} as follows.
Set $\Gamma=\pi_1(M)$ with holonomy
$\rho:\pi_1(M)\to\mathrm{PSL}(2,\mathbb{C})$.
For a loxodromic element $\gamma$, the boundary direction
$\xi(\gamma)=\hat{n}(\gamma)\in S^2$ is the geodesic axis extracted
via the Pauli decomposition of $\log\rho(\gamma)$,
and the inner product becomes the Gaussian kernel
$\langle\xi(f)|\xi(g)\rangle=\exp(-\theta_{fg}^2/(2\sigma^2))$,
$\theta_{fg}=\arccos|\hat{n}(f)\cdot\hat{n}(g)|$.

For $M_{\mathrm{CKM}}=m006(-5,2)$ with word triple
$\{\mathrm{aaB},\mathrm{AbA},\mathrm{AAb}\}$:
the homological class collision
$[\mathrm{AbA}]=[\mathrm{AAb}]=2$ in $H_1(M)=\mathbb{Z}/5$
forces $J_{\mathrm{CP}}=0$ via the No Continuous Drift Corollary,
consistent with $|J_{\mathrm{CKM}}|\approx3\times10^{-5}$;
the Frobenius fitness is $0.016482$ against PDG~2024.

For $M_{\mathrm{PMNS}}=m003(-2,3)$ with Borel construction on
$\{\mathrm{aa},\mathrm{aaB},\mathrm{baa}\}$:
the distinct classes $[\mathrm{aa}]=1$, $[\mathrm{aaB}]=4$,
$[\mathrm{baa}]=3$ (mod~5) ensure $J_{\mathrm{CP}}\neq0$,
consistent with large leptonic CP violation;
the fitness is $0.005087$, the global minimum of the construction.

\section{Conclusions}
\label{sec:conclusion}

Under the discrete boundary assignment, the set of admissible mixing
matrices is finite modulo rephasings (Theorem~\ref{thm:discreteness}),
and bounded deformations cannot produce continuous drift in physical
observables (Theorem~\ref{thm:stability}).
This provides a model-independent kinematic explanation for why the
CKM and PMNS mixing matrices cannot vary continuously: they are
selected from a discrete geometric set determined by the topology
of the underlying manifold.

\section*{Data Availability}
No research data generated.
Companion scripts: \url{https://github.com/drmlgentry/hyperbolic-flavor-scan}.

\section*{Acknowledgments}
The author thanks the developers of SnapPy and SageMath.

\begin{thebibliography}{99}

\bibitem{PDG}
S.~Navas et al.\ (PDG),
Phys.\ Rev.\ D \textbf{110}, 030001 (2024).

\bibitem{AltarelliFeruglio}
G.~Altarelli, F.~Feruglio,
Rev.\ Mod.\ Phys.\ \textbf{82}, 2701 (2010).

\bibitem{KingLuhn}
S.~F.~King, C.~Luhn,
Rept.\ Prog.\ Phys.\ \textbf{76}, 056201 (2013).

\bibitem{Jarlskog}
C.~Jarlskog,
Z.\ Phys.\ C \textbf{29}, 491 (1985).

\bibitem{Ratcliffe}
J.~G.~Ratcliffe,
\textit{Foundations of Hyperbolic Manifolds}, Springer (2006).

\bibitem{BridsonHaefliger}
M.~R.~Bridson, A.~Haefliger,
\textit{Metric Spaces of Non-Positive Curvature}, Springer (1999).

\bibitem{Gentry:CKM}
M.~L.~Gentry,
\textit{Quark Mixing from Hyperbolic Geometry},
SSRN \href{https://doi.org/10.2139/ssrn.6583550}{6583550} (2026);
J.\ Geom.\ Phys.\ JGP13076.

\bibitem{Gentry:PMNS}
M.~L.~Gentry,
\textit{Lepton Mixing from Borel Structure of Hyperbolic Holonomy},
SSRN \href{https://doi.org/10.2139/ssrn.6583549}{6583549} (2026);
Ann.\ Henri Poincar\'e AHPO-D-26-00255.

\bibitem{Gentry:Unified}
M.~L.~Gentry,
\textit{Hyperbolic Flavor Geometry},
SSRN \href{https://doi.org/10.2139/ssrn.6775158}{6775158} (2026);
Phys.\ Rev.\ D DS14327.

\end{thebibliography}

\end{document}